\secn{Supplemental Problem Bank}
All by hand (no calculator, no Desmos). Questions only; no answers. ``Show/justify'' = full proof or working expected. $\mathbb F$ is an arbitrary field unless stated.

\sub{1. Row Reduction}
\bi{RR1}{Which of $A=\begin{pmatrix}1&0&2\\0&1&-1\\0&0&0\end{pmatrix}$, $B=\begin{pmatrix}1&0&2\\0&1&0\\0&0&1\end{pmatrix}$, $C=\begin{pmatrix}1&2&0\\0&1&0\\0&0&0\end{pmatrix}$ are in RREF? For each other matrix, name the violated condition.}
\bi{RR2}{(a) Write an augmented matrix in REF of an inconsistent system. (b) Write one with a unique solution and one with infinitely many. (c) For $\left[\begin{smallmatrix}1&2&-1&0&5\\0&1&3&4&-2\\0&0&0&1&7\end{smallmatrix}\right]$ state the basic and free variables.}
\bi{RR3}{\textit{(challenge)} Let $S'$ arise from $S$ by one elementary row operation ($R_i\leftrightarrow R_j$, $cR_i$, $R_i+\lambda R_j$). Prove $x$ solves $S\iff x$ solves $S'$; conclude $S$ and $S_{\mathrm{rref}}$ have equal solution sets.}
\bi{RR4}{Reduce to RREF and give LI spanning sets for $\mathrm{col}(A)$ and $\mathrm{row}(A)$: (a) $\begin{pmatrix}1&2&-1\\2&4&-2\\1&1&0\end{pmatrix}$, (b) $\begin{pmatrix}2&1&0&3\\1&0&1&2\\3&1&1&5\end{pmatrix}$, (c) $\begin{pmatrix}1&i&0\\2&2i&1\\1&i&1\end{pmatrix}$ over $\mathbb C$.}
\bi{RR5}{Let the homogeneous system for $a_1v_1+\dots+a_nv_n=0$ reduce to REF with basic indices $B$, free indices $F$. Prove (a) $v_f\in\mathrm{span}\{v_b:b\in B\}$ for $f\in F$; (b) $\mathrm{span}\{v_b\}=\mathrm{span}(S)$; (c) $\{v_b:b\in B\}$ is LI. Deduce $\mathrm{col}(A)=\mathrm{span}\{c_i: i\text{ pivot column of }\hat A\}$.}

\sub{2. Linear Systems}
\bi{LS1}{For each system, (i) Gaussian elimination to REF, (ii) Gauss--Jordan to RREF, (iii) basic/free variables, (iv) for the homogeneous version a spanning set of the solution space, (v) for the non-homogeneous version the solution set as $p+\mathrm{span}(\dots)$: \\ (a) $x{+}y{+}z=2,\ 2x{+}2y{+}2z=4,\ x{-}y{+}z=0$; (b) $x{+}y{+}z=3,\ 2x{-}y{+}z=1,\ -x{+}y{+}2z=4$; (c) $x{+}y=2,\ 2x{-}y=1,\ 3x=3$; (d) $x{+}y{+}z=1,\ 2x{+}2y{+}2z=2,\ x{+}y{+}z=3$.}
\bi{LS2}{\textit{Set-up only.} Let $v_1=(1,i,0,1)$, $v_2=(2,1,i,0)$, $v_3=(3,1{+}i,i,1)\in\mathbb C^4$. (a) Using an iff chain, construct a system whose solutions decide if $\{v_1,v_2,v_3\}$ is LI. (b) Likewise decide if $w=(4,2{+}i,i,1)\in\mathrm{span}\{v_1,v_2,v_3\}$.}
\bi{LS3}{(a) $W=\{(a,b,c,d)\in\mathbb R^4:a+b+c+d=0\}$: build a system, solve by Gauss--Jordan, write $W=\mathrm{span}(S)$ with $S$ LI. (b) $W=\{(x,y,z):2x+y=z\}$: construct $S$ and prove $\mathrm{span}(S)=W$ by double inclusion.}
\bi{LS4}{Let $x,y$ solve $a_{r1}v_1+\dots+a_{rn}v_n=0$ $(r=1..k)$. (a) Show $cx+y$ solves a generic row. (b) Deduce the solution set is a subspace of $\mathbb F^n$.}
\bi{LS5}{Solve $Ax=b$ in $\mathbb C^3$: (a) $A=\begin{pmatrix}1&1&0\\0&1-i&1\\0&0&i\end{pmatrix}$, $b=(i,1,1{+}i)^T$; (b) $A=\begin{pmatrix}1&0&1\\0&1&i\\0&0&0\end{pmatrix}$, $b=(1{+}i,2,0)^T$; (c) $A=\begin{pmatrix}1{+}i&1&0\\0&1&i\\0&0&1\end{pmatrix}$, $b=(2,1{+}i,-i)^T$.}

\sub{3. Matrix Operations (including inverses)}
\bi{MO1}{For $A,B\in M_{n\times k}(\mathbb C)$ expand with no brackets: (a) $(2A-3B)^*$, (b) $((3-i)A+2B)^*$, (c) $3iA(iB-4C)^*$. (d) Compute $(2A-3B)^*$ for $A=\begin{pmatrix}1&2\\-i&3\end{pmatrix}$, $B=\begin{pmatrix}0&-1\\2&i\end{pmatrix}$.}
\bi{MO2}{Prove for $X,Y\in M_{m\times n}(\mathbb C)$, $c\in\mathbb C$: $(X{+}Y)^T$, $(X{+}Y)^*$, $(cX)^*=\bar cX^*$, and for $A\in M_{n\times l}$, $B\in M_{l\times k}$: $(AB)^*=B^*A^*$. Verify the first three for $A=\begin{pmatrix}1{+}i&2\\-3i&4\end{pmatrix}$, $B=\begin{pmatrix}2&-1{+}i\\5&3i\end{pmatrix}$, $c=2-i$.}
\bi{MO3}{Compute: (a) $\begin{pmatrix}1{+}i&2\\0&1{-}i\end{pmatrix}\begin{pmatrix}i&1\\2&-i\end{pmatrix}$; (b) $(AB)^*$ for $A=\begin{pmatrix}3i&4\\2{-}i&2\end{pmatrix}$, $B=\begin{pmatrix}1&-i\\0&2i\end{pmatrix}$; (c) $\begin{pmatrix}1{+}i&2&0\\-i&1&3\\2&i&1{-}i\end{pmatrix}\begin{pmatrix}2{-}i&1&i\\1{+}i&0&2\\i&1{-}i&1\end{pmatrix}$.}
\bi{MO4}{\textit{(enrichment)} $\begin{pmatrix}1{+}i&2&0&-i\\0&1{-}i&3i&2\\2&0&i&1\\-i&1&0&2{+}i\end{pmatrix}\begin{pmatrix}1\\i\\-1\\2\end{pmatrix}$.}
\bi{MO5}{For $A,B\in M$ with compatible shapes prove entrywise (write $(Au)_{i1}=\sum_ja_{ij}u_{j1}$): $A(u{+}v)=Au+Av$; $(cA)u=c(Au)=A(cu)$; $A0=0$; $A(cu{+}dv)=cAu+dAv$; $I_nu=u$; $(A{+}B)u=Au+Bu$; $A(B{+}C)=AB+AC$.}
\bi{MO6}{Construct $A,B\in M_2(\mathbb R)$, both non-invertible, with $A+B$ invertible.}
\bi{MO7}{Is it invertible? (a) $\begin{pmatrix}1&2\\3&4\end{pmatrix}$, (b) $\begin{pmatrix}1&2\\2&4\end{pmatrix}$, (c) $\begin{pmatrix}1&0&2\\0&1&3\\0&0&1\end{pmatrix}$, (d) $\begin{pmatrix}1&2&3\\2&4&6\\1&0&1\end{pmatrix}$, (e) $\begin{pmatrix}1{+}i&2\\3&1{-}i\end{pmatrix}$.}
\bi{MO8}{\textit{(enrichment)} $f(x)=x^2+3x$, $g(x)=2x+x^2$. Compute $f(2x+3)$, $(2f+3)(x)$, $(2f+2g)(x)$, $(fg)(x)$, $(f\circ g)(x)$. State when juxtaposition means product and when composition.}
\bi{INV1}{For each $T$: choose ordered bases, compute the matrix, use Gauss--Jordan to decide invertibility, and if invertible give $T^{-1}$ explicitly; else say why not. (a) $T(x,y)=(2x{+}y,x{-}y)$; (b) $p\mapsto p'$ on $\mathbf P_2(\mathbb R)$; (c) $A\mapsto A+A^T$ on $M_2(\mathbb R)$; (d) $T(x,y,z)=(x{+}y,y{+}z,x{+}z)$; (e) $a{+}bx{+}cx^2\mapsto a{+}2bx{+}3cx^2$; (f) $p\mapsto\int_0^xp$, $\mathbf P_2\to\mathbf P_3$; (g) $A\mapsto\begin{pmatrix}1&1\\0&1\end{pmatrix}A$; (h) $(x,y,z)\mapsto(x,y,0)$; (i) $a{+}bx{+}cx^2\mapsto a+(b{+}a)x+cx^2$; (j) $T:M_2(\mathbb C)\to\mathbf P_3(\mathbb C)$, $\left[\begin{smallmatrix}a&b\\c&d\end{smallmatrix}\right]\mapsto(a{+}bi)+(b{+}ci)x+(c{+}di)x^2+dx^3$ (find $([T]_\alpha^\beta)^{-1}$ and $T^{-1}$).}
\bi{INV2}{(a) If $AB=I_n$ in $M_{n\times n}(\mathbb F)$, explain via linear maps why $BA=I_n$. (b) Prove an inverse of a linear map is unique; same for matrices. (c) Prove $T^{-1}$ is linear when $T$ is an invertible linear map.}
\bi{INV3}{$T(x,y,z)=(x{+}y,y{+}z,x{+}z)$ on $\mathbb F^3$. Without computing $\ker$ or image, decide injectivity and surjectivity; how does the answer depend on $\mathbb F$? (Hint: $\det=2$.)}

\sub{4. Determinants}
\bi{DT1}{Using only multilinearity, alternation and $\det(e_1,\dots,e_n)=1$ (no $ad-bc$, no cofactors) compute $\det$ of (a) $\begin{pmatrix}1&2\\3&4\end{pmatrix}$, (b) $\begin{pmatrix}2&3\\2&-4\end{pmatrix}$, (c) $\begin{pmatrix}1&i&0\\2&3&0\\0&0&2{-}i\end{pmatrix}$, (d) $\begin{pmatrix}1{+}i&1&0\\0&1&i\\0&0&2\end{pmatrix}$, (e) $\begin{pmatrix}1&2&3\\2&4&6\\1&0&1\end{pmatrix}$.}
\bi{DT2}{Prove (induction where indicated): $\det A^T=\det A$; $\det\bar A=\overline{\det A}$ (challenge) hence $\det A^*=\overline{\det A}$; triangular $\Rightarrow\det A=\prod a_{ii}$; $\det kA=k^n\det A$; equal rows $\Rightarrow0$; swapping rows negates; adding a multiple of a row leaves $\det$ unchanged.}
\bi{DT3}{Using only the multilinear and alternating properties prove: if $B$ is $A$ with row $i$ multiplied by $k$, then $\det B=k\det A$.}
\bi{DT4}{Cofactor expansion: (a) $\begin{pmatrix}2&1&0\\-1&3&4\\0&2&5\end{pmatrix}$, (b) $\begin{pmatrix}1&0&2\\3&1&4\\2&0&1\end{pmatrix}$, (c) $\begin{pmatrix}1{+}i&1&2\\2&1{+}i&1\\1&2&1{+}i\end{pmatrix}$, (d) $\begin{pmatrix}1&2&0&1\\0&1&3&2\\0&0&2&1\\0&0&0&4\end{pmatrix}$.}
\bi{DT5}{(a) Area of the triangle $(1,2),(4,1),(6,5)$; of $(-1,3),(2,0),(5,4)$. (b) Area of the parallelogram spanned by $(2,-1),(4,3)$. (c) Volume of the parallelepiped spanned by $(1,0,2),(2,1,1),(3,-1,4)$. (d) Volume of the tetrahedron $(0,0,0),(1,2,0),(3,1,1),(2,0,4)$. (e) Explain the factors $\tfrac12$ and $\tfrac16$.}
\bi{DT6}{Decide whether Cramer's rule applies (show $\det\ne0$); if so solve: (a) $x{+}2y=5,\ 3x{+}4y=11$; (b) $2x{-}y=1,\ 4x{-}2y=2$; (c) $x{+}y{+}z=6,\ 2x{+}3y{+}z=10,\ x{+}4y{+}5z=20$; (d) $x{+}2y{+}z=4,\ 2x{+}4y{+}2z=8,\ x{+}y{+}z=3$; (e) $3x{-}y{+}2z=7,\ 2x{+}y{-}z=1,\ x{+}4y{+}z=9$; (f) $x{+}y=2,\ 2x{+}2y=4,\ 3x{+}3y=6$; (g) $2x{+}y{-}z=3,\ x{-}y{+}2z=0,\ 3x{+}2y{+}z=5$; (h) $x{+}y=3,\ 2y{+}z=4,\ x{+}z=5$.}
\bi{DT7}{State the characterization of $\det:M_{n\times n}(\mathbb F)\to\mathbb F$ by its multilinear and alternating properties.}

\sub{5. Vector Spaces}
\bi{VS1}{Prove for a field: additive identity unique; additive inverse unique; $0\cdot c=0$; $(-1)c=-c$. Prove for a vector space $V$: zero vector unique; inverse unique; cancellation ($x{+}y=x{+}z\Rightarrow y=z$); $0v=c0=0$; $(-1)x=-x$. List two such properties that are not axioms.}
\bi{VS2}{$\boxplus:S\times S\to S$ is commutative and has $0$ with $0\boxplus s=s$ for all $s$. Prove $0$ is the unique such element.}
\bi{VS3}{Prove: (a) $c(x+y)=cx+cy$ in $\mathbb F^n$; (b) $(c+d)f=cf+df$ in $\mathcal F(\mathbb C)$; (c) $M_{n\times k}(\mathbb C)$ and $\mathcal F(\mathbb C)$ are vector spaces; (d) $\mathcal F_{\mathrm{even}}(\mathbb R)$ is a vector space. Structure: given / to prove / technique.}
\bi{VS4}{\textit{(challenge)} On $\mathbb R^+=\{x>0\}$ set $x\boxplus y=xy$, $c\boxdot x=x^c$. Show $(\mathbb R^+,\boxdot,\boxplus)$ is an $\mathbb R$-vector space; identify its zero and additive inverses.}
\bi{VS5}{(a) Prove $\mathbb C$ is closed under $+$ and $\times$ and has an additive identity. (b) Compute $(2-3i)(1+2i)$; $1/i$; $(1+i)^{-1}$; $\frac{3+2i}{1-i}$; $-(2-5i)$; $(2+i)^{-1}$.}
\bi{VS6}{Subspace of $M_{m\times n}(\mathbb F)$ (proof or counterexample): $A_{11}=0$; $A_{11}=1$; $A_{12}=2A_{21}$; symmetric; skew-symmetric; diagonal. Subspace of $\mathcal F(\mathbb C)$: $f(0)=0$; $f(0)=1$; even; odd; real-valued.}
\bi{VS7}{Subspace? If yes prove it; if not give a counterexample: (a) $\{(x,y,z):x{+}y{-}z=0\}$; (b) $\{A\in M_2(\mathbb C):A=A^*\}$; (c) $\{p\in\mathbf P_2(\mathbb R):p(1)=1\}$; (d) $\{A\in M_2(\mathbb C):\mathrm{tr}A=0\}$; (e) $\{p\in\mathbf P_3(\mathbb R):p(0)+p(1)=0\}$; (f) $\{x{+}y{-}z=2\}$; (g) $\{f\in\mathbf P_2(\mathbb C):f(1)=f(2),\,f(3)=1\}$; (h) $\{A\in M_2(\mathbb R):\det A=0\}$; (i) $\{f\in\mathbf P_3(\mathbb R):f(0)=f(1),\,f(2)=0\}$.}
\bi{VS8}{$W_1,W_2\le V$. Prove $W_1\cap W_2\le V$. Is $W_1\cup W_2$ a subspace? When?}
\bi{VS9}{Prove: $S\subseteq\mathrm{span}(S)$; $S\subseteq T\Rightarrow\mathrm{span}(S)\subseteq\mathrm{span}(T)$; $\mathrm{span}(\mathrm{span}S)=\mathrm{span}S$; $\mathrm{span}(S)$ is a subspace (handle $S=\varnothing$).}
\bi{VS11}{True, false, or sometimes true (justify): $S\subseteq T$ and (a) $S$ LI $\Rightarrow T$ LI; (b) $S$ LD $\Rightarrow T$ LD; (c) $T$ LI $\Rightarrow S$ LI; (d) $T$ LD $\Rightarrow S$ LD.}
\bi{VS12}{(a) Define a subspace of $M_2(\mathbb C)$ of dimension 2. (b) Give $S\subseteq\mathbf P_3(\mathbb R)$ with $\mathrm{span}S=\mathbf P_3(\mathbb R)$ and $S$ LD. (c) $S=\{x_1..x_n\}$ spans $V$ and is LD: what can be said about $\dim V$?}
\bi{VS13}{True or false with proof/counterexample: (a) $W\le V$, $V$ infinite-dim.\ $\Rightarrow W$ infinite-dim.; (b) $\Rightarrow W$ finite-dim.; (c) $V$ finite-dim.\ $\Rightarrow W$ infinite-dim.; (d) $\Rightarrow W$ finite-dim. (e) Prove $\mathbf P(\mathbb R)=\bigcup_n\mathbf P_n(\mathbb R)$ is not finite-dimensional.}
\bi{VS14}{(a) $W=\{A\in M_2(\mathbb F):A^T=A\}$: basis; if $W\cong\mathbb F^n$ find $n$. (b) Explain why $\dim\mathbb F^n,M_{n\times k}(\mathbb F),\mathbf P_n(\mathbb F)$ do not depend on $\mathbb F$.}

\sub{6. Linear Independence}
\bi{LI1}{LI or LD (use an iff setup): (a) $\{(1,i,0),(2,1,i),(1,0,1)\}\subseteq\mathbb C^3$; (b) $\left\{\left[\begin{smallmatrix}1&i\\0&1\end{smallmatrix}\right],\left[\begin{smallmatrix}i&1\\0&i\end{smallmatrix}\right],\left[\begin{smallmatrix}1{+}i&2\\0&1{+}i\end{smallmatrix}\right]\right\}\subseteq M_2(\mathbb C)$; (c) $\{1{+}x,\,2{-}x{+}x^2,\,1{+}3x{+}x^2\}\subseteq\mathbf P_2(\mathbb R)$; (d) $\left\{\left[\begin{smallmatrix}1&i\\0&1\end{smallmatrix}\right],\left[\begin{smallmatrix}2&2i\\0&2\end{smallmatrix}\right],\left[\begin{smallmatrix}1&0\\0&0\end{smallmatrix}\right]\right\}$; (e) $\{I,2I,3I\}\subseteq M_2(\mathbb R)$ (set-up only).}
\bi{LI2}{Is $x$ in the span? (a) $(3,4,1)\in\mathrm{span}\{(1,1,0),(2,1,1),(1,2,0)\}$; (b) $(1{+}i,2)\in\mathrm{span}\{(1,i),(2,1)\}\subseteq\mathbb C^2$; (c) $\left[\begin{smallmatrix}1{+}i&2\\0&1{+}i\end{smallmatrix}\right]\in\mathrm{span}\left\{\left[\begin{smallmatrix}1&i\\0&1\end{smallmatrix}\right],\left[\begin{smallmatrix}i&1\\0&i\end{smallmatrix}\right]\right\}$; (d) $1{+}2x{+}x^2\in\mathrm{span}\{1{+}x,1{-}x{+}x^2\}$; (e) $1{+}x{+}x^2\in\mathrm{span}\{1{+}x,x^2,x^3\}$; (f) $2x{+}1\in\mathrm{span}\{x^2{+}3x{+}2,x{+}1,x^2{+}2\}$ (set-up); (g) $x^2{-}x\in\mathrm{span}\{x^2{+}1,x^2{+}x{+}1\}$ (system, RREF, conclude).}
\bi{LI3}{$V$ complex; $u,v,w\in V$. Prove $\{u,v,w\}$ LI $\iff\{u{+}v,u{+}w,v{+}w\}$ LI. (Where does $2\ne0$ matter?)}
\bi{LI4}{$S$ LI, $S'\subsetneq S$, $W=\mathrm{span}(S)$. Prove $\mathrm{span}(S')\ne W$.}
\bi{LI5}{$S=\{(1,2,3),(3,2,1),(1,1,1)\}$. (a) Write the system for LI and justify it. (b) Gaussian elimination to REF. (c) Decide LI/LD and, if LD, a LI $S'\subseteq S$ that is a basis of $\mathrm{span}S$.}
\bi{LI6}{$T\in\mathcal L(V,W)$. (a) $\{Tv_1..Tv_n\}$ LI $\Rightarrow\{v_1..v_n\}$ LI. (b) $T$ injective and $\{v_i\}$ LI $\Rightarrow\{Tv_i\}$ LI. (c) $T$ preserves linear dependence.}

\sub{7. Basis and Dimension}
\bi{BD1}{Prove $\{e_i\}$ is a basis of $\mathbb F^n$, $\{E_{ij}\}$ of $M_{n\times k}(\mathbb F)$, $\{1,x,\dots,x^n\}$ of $\mathbf P_n(\mathbb F)$; state the dimensions.}
\bi{BD2}{Basis and dimension: (a) $\{A\in M_2(\mathbb R):\mathrm{tr}A=0\}$; (b) $\{(x,y,z):x=2y\}$; (c) $\{f\in\mathbf P_3(\mathbb R):f(0)=f(1)=0\}$; (d) $\{A\in M_2(\mathbb C):A_{21}=0\}$; (e) $\{A\in M_2(\mathbb R):a{+}d=0,\ b=c\}$; (f) $\{p\in\mathbf P_3(\mathbb R):p(0)=p(1),\,p(2)=0\}$; (g) $\{(x,y,z,w):x{+}y{-}z=0,\ y{+}w=0\}$; (h) $\{A\in M_2(\mathbb R):\mathrm{tr}A=0,A=A^T\}$; (i) $\{A\in M_3(\mathbb C):A=A^*\}$ as a real vector space (why is it not a complex subspace?).}
\bi{BD3}{Counting arguments: $\{(1,0,0),(0,1,0),(0,0,1),(1,1,1)\}$ not LI; $\{(1,0,0),(0,1,0)\}$ not spanning $\mathbb R^3$; five matrices in $M_2(\mathbb R)$ are LD; five matrices in $M_{3\times2}(\mathbb R)$ do not span; $\{1,x,x^2,x^3\}\subseteq\mathbf P_2$ LD; $\{1,x\}$ does not span $\mathbf P_2$; ten matrices in $M_3(\mathbb R)$ are not a basis; $n$ polynomials in $\mathbf P_n(\mathbb R)$ are not a basis.}
\bi{BD4}{Extend the LI set $\left\{\left[\begin{smallmatrix}1&0\\0&0\end{smallmatrix}\right],\left[\begin{smallmatrix}0&1\\1&0\end{smallmatrix}\right],\left[\begin{smallmatrix}0&0\\1&1\end{smallmatrix}\right]\right\}$ to a basis of $M_2(\mathbb R)$ using standard basis vectors; justify the choice (not every $E_{ij}$ works).}
\bi{BD5}{$\dim V=n$, $S$ LI with $|S|=n$. Using the Extension Lemma and Dimension Bound Theorem, prove $S$ is a basis.}
\bi{BD6}{Find an LI subset with the same span: (a) $\{(1,2,1),(2,4,2),(1,1,0)\}$; (b) $\left\{\left[\begin{smallmatrix}1&2\\0&1\end{smallmatrix}\right],\left[\begin{smallmatrix}2&4\\0&2\end{smallmatrix}\right],\left[\begin{smallmatrix}1&1\\0&0\end{smallmatrix}\right]\right\}$; (c) $\{1{+}x,2{+}2x,x\}$.}

\sub{8. Linear Transformations}
\bi{LT1}{Linear? Prove or disprove: $T(x,y)=(2x{-}y,x{+}3y)$; $T(A)=A+A^T$; $T(x,y,z)=x-2y+z$; $T(x)=x^2$; $T(p)=p(1)$; $T(p)=p'$; $T(x,y)=(0,x{+}y)$; $T(A)=\mathrm{tr}A+1$ ($M_2\to\mathbb R$); $T(A)=A+I$ on $M_2(\mathbb R)$.}
\bi{LT2}{Prove for linear $T$: $T(0)=0$; $T(-v)=-T(v)$; $T(\sum c_iv_i)=\sum c_iT(v_i)$; LD preserved; $T(U)$ is a subspace for $U\le V$; $\mathrm{image}(T)\le W$.}
\bi{LT3}{(i) $[T]_\alpha^\beta$, (ii) $[v]_\alpha$, (iii) $[T(v)]_\beta$, (iv) verify $[T(v)]_\beta=[T]_\alpha^\beta[v]_\alpha$: (a) $T(x,y)=(2x{+}y,x{-}y)$, $\alpha$ standard, $\beta=\{(1,1),(1,-1)\}$, $v=(3,1)$; (b) $p\mapsto p'$ on $\mathbf P_2$, $\alpha=\{1,x,x^2\}$, $\beta=\{1,x,1{+}x^2\}$, $p=2{+}3x{+}x^2$; (c) $T(A)=A+A^T$ on $M_2$ with $\beta=\{E_{11},E_{12}{+}E_{21},E_{22},E_{12}{-}E_{21}\}$, $A=\left[\begin{smallmatrix}1&2\\3&4\end{smallmatrix}\right]$; (d) $T(a{+}bx{+}cx^2)=\left[\begin{smallmatrix}a&b\\c&a\end{smallmatrix}\right]$ and $T(a,b,c)=\left[\begin{smallmatrix}a{+}b&a{-}c\\b&a{+}b\end{smallmatrix}\right]$ in standard bases.}
\bi{LT4}{$M=\begin{pmatrix}1&2&0\\0&-1&3\\2&0&1\end{pmatrix}$, $\alpha$ standard of $\mathbb R^3$, $\beta=\{1,x,x^2\}$. Define $T_M(a,b,c)$ explicitly.}
\bi{LT5}{$T:\mathbb R^3\to\mathbf P_2$ linear with $T(1,2,3)=2+3x$, $T(1,1,1)=x^2$. Find $T(0,1,2)$.}
\bi{LT6}{(a) $I$ and $Z$ are linear on $V$; $I\circ T=T\circ I=T$; $T\circ Z=Z$. (b) $(Tp)(x)=\int_0^xp(t)dt$, $\mathbf P_n\to\mathbf P_{n+1}$: linear; $T(x^k)$; $(Tp)(0)=0$.}
\bi{LT8}{(a) Prove for $A=[T]_\alpha^\beta$: $(b_1,\dots,b_k)$ solves $Ax=0\iff b\in\mathrm{null}(A)\iff\sum b_iv_i\in\ker T$; and $w\in\mathrm{image}(T)\iff[w]_\beta\in\mathrm{col}(A)$. (b) Bases of $\ker T$, $\mathrm{image}T$ and rank--nullity for $T(x,y)=(2x{+}y,x{-}y)$; $p\mapsto p'$; $A\mapsto A+A^T$. (c) $T:\mathbf P_2(\mathbb C)\to\mathbb C^3$, $a{+}bx{+}cx^2\mapsto(a{+}ib{+}2c,\ (1{+}2i)a{+}b{+}(1{+}i)c,\ 2ia{+}(2{-}i)b{+}c)$: matrix, null-space basis, $\ker T$, col-space basis, $\mathrm{image}T$, rank--nullity. (d) $A=\left[\begin{smallmatrix}1&-2&0\\3&1&4\end{smallmatrix}\right]$, $\alpha=$ std of $\mathbb R^2$, $\beta=\{1,x,x^2\}$: define $T$, find $\ker,\mathrm{image}$.}
\bi{LT9}{$T\in\mathcal L(V)$. Can $\mathrm{image}T=\ker T$ for $V=\mathbb R^2$? $\mathbb R^3$? State the condition on $\dim V$; prove existence when it holds and impossibility otherwise.}
\bi{LT10}{(a) $T\in\mathcal L(V,W)$, $V$ finite-dim.: $\ker T=\ker T^2\iff\mathrm{image}T=\mathrm{image}T^2$. (b) $V'\le V$, $W'\le W$: $T(V')$ and $T^{-1}(W')$ are subspaces. (c) Choose $W'$ with $T^{-1}(W')=\ker T$ and $V'$ with $T(V')=\mathrm{image}T$; state the most general such conditions.}
\bi{LT11}{Construct: (a) linear $T:\mathbb R^3\to\mathbb R^3$ with $\ker T=\{0\}$; (b) linear $T:\mathbf P_2(\mathbb R)\to\mathbb R^3$ that is not an isomorphism; (c) state $\mathrm{image}(P_{(1,2,1)})$.}
\bi{LT12}{Yes/No with a one-line dimension argument: $T:\mathbb R^2\to\mathbb R^3$ surjective? $\mathbb R^3\to\mathbb R^2$ injective? $\mathbf P_1\to\mathbb R^2$ surjective? $\mathbb R^3\to\mathbb R^3$ injective? $M_{3\times2}(\mathbb R)\to\mathbf P_6(\mathbb R)$ surjective? $V\to M_{2\times4}(\mathbb R)$ surjective: what for $\dim V$ (and if $V$ infinite)? Rank: $\mathbb R^5\to\mathbb R^3$ with nullity 4; $M_2\to\mathbb R^4$ with nullity 3; $\mathbf P_2\to\mathbb R^3$ with nullity 0. Compute $\dim\mathcal L(\mathbb R^5,\mathbb R^3)$, $\dim\mathcal L(\mathbf P_2,\mathbb R^3)$, $\dim\mathcal L(M_2(\mathbb R),\mathbb R^4)$, $\dim\mathcal L(M_{2\times4}(\mathbb C),\mathbf P_2(\mathbb C))$; $\Phi\in\mathcal L(\mathcal L(\mathbb R^2,\mathbb R^3),\mathbb R^7)$: can $\Phi$ be injective? surjective?}
\bi{LT13}{Prove for $T\in\mathcal L(V,W)$: injective $\Rightarrow$ LI preserved; surjective $\Rightarrow$ spanning sets preserved; ($\dim V<\infty$) injective $\iff\dim\mathrm{image}T=\dim V$; ($\dim W<\infty$) surjective $\iff\dim\mathrm{image}T=\dim W$. State the equivalent conditions for injectivity. On $\mathbb F^{\mathbb N}$: right shift is not surjective; left shift is not injective.}
\bi{LT14}{(a) $S,T\in\mathcal L(V,W)$: $S{+}T,\ cT\in\mathcal L(V,W)$. (b) $R,S,T\in\mathcal L(V)$: $T(R{+}S)=TR+TS$. (c) $R\in\mathcal L(U,V),S\in\mathcal L(V,W),T\in\mathcal L(W,X)$: $SR$ linear and $T(SR)=(TS)R$. (d) $\mathcal L(V,W)$ is a vector space.}
\bi{LT15}{$T:U\to V$, $S:V\to W$. (a) $ST$ injective $\Rightarrow T$ injective; $ST$ surjective $\Rightarrow S$ surjective. (b) Give $U,V,W,T,S$ with $ST$ an isomorphism but $S$ not injective and $T$ not surjective. (c) $ST$ injective: $S$ injective $\iff\dim U=\dim V=\dim W$. (d) $ST$ iso $\Rightarrow\dim U\ge\dim V\ge\dim W$.}

\sub{9. Eigenvalues and Eigenvectors}
\bi{EV3}{Eigenvalues and eigenspace bases: (a) $T(x,y)=(3x,2y)$; (b) $T(x,y)=(x{+}y,y)$ (is it diagonalizable?); (c) $p\mapsto p'$ on $\mathbf P_2(\mathbb R)$ (matrix in $\{1,x,x^2\}$); (d) $A\mapsto A^T$ on $M_2(\mathbb R)$ (describe eigenspaces).}
\bi{EV4}{(a) Eigenvalues of $\begin{pmatrix}2{+}i&3&0\\0&-1{+}2i&4\\0&0&5{-}i\end{pmatrix}$ and an eigenvector for $2{+}i$. (b) Eigenvalues and eigenspace bases of $A=\begin{pmatrix}1{+}i&1&0\\0&1&i\\0&0&2\end{pmatrix}$. (c) Define eigenvalue of a linear transformation.}
\bi{EV5}{(a) Define $T:\mathbb R^2\to\mathbb R^2$ with exactly one eigenvalue and give its matrix. (b) $T$ invertible: what is $E_0$? (c) Eigenvalues of $T(x,y,z)=(x{+}y{+}z,\ 2y{-}3z,\ 4z)$. (d) For $T(x,y,z)=(2x,y,-2z)$: diagonalizable? invertible? Why?}
\bi{EV6}{(a) Prove $\overline{zw}=\bar z\bar w$. (b) Prove $\overline{Av}=\bar A\bar v$. (c) $A\in M_n(\mathbb R)$, $Av=\lambda v$, $v\ne0$: show $A\bar v=\bar\lambda\bar v$ (and $\bar v\ne0$). (d) $\lambda_1\ne\lambda_2$ eigenvalues of $A$: prove $E_{\lambda_1}\cap E_{\lambda_2}=\{0\}$.}
\bi{EV8}{$A=\begin{pmatrix}1&-1&2\\1&0&-5\\0&1&3\end{pmatrix}$: compute $p(\lambda)=\det(A-\lambda I)$ and all eigenvalues (real and complex).}
\bi{EV9}{$P$ = orthogonal projection of $\mathbb R^2$ onto $\mathrm{span}\{(1,2)\}$. Find $P(e_1),P(e_2)$, $[P]$, eigenvalues, eigenspace bases, and show the eigenspaces are orthogonal. In general, for $u=(a,b)\ne0$ show $u$ has eigenvalue 1 and $(-b,a)$ eigenvalue 0.}

\sub{10. Diagonalization (including change of basis)}
\bi{EV1}{(a) $A=MBM^{-1}$, $p(x)=\sum a_ix^i$: prove $A^r=MB^rM^{-1}$ and $p(A)=Mp(B)M^{-1}$; $p(\mathrm{diag}(d_i))=\mathrm{diag}(p(d_i))$; $[T]_\alpha^\alpha=M D M^{-1}$ $\Rightarrow p([T]_\alpha^\alpha)=Mp(D)M^{-1}$. (b) $p(D)$: $D=\mathrm{diag}(2,3)$, $x^2{-}4x{+}1$; $D=\mathrm{diag}(1,-2,4)$, $x^3{-}2x{+}5$; $D=\mathrm{diag}(0,1,-1)$, $2x^3{+}x^2{-}1$. (c) $p(A)$ with $A=MDM^{-1}$: $M=\left[\begin{smallmatrix}1&1\\0&1\end{smallmatrix}\right]$, $D=\mathrm{diag}(2,5)$, $x^2{+}1$; $M=\left[\begin{smallmatrix}1&0\\2&1\end{smallmatrix}\right]$, $D=\mathrm{diag}(-1,3)$, $x^2{-}2x$; $M=\left[\begin{smallmatrix}1&0&0\\0&1&1\\0&0&1\end{smallmatrix}\right]$, $D=\mathrm{diag}(1,2,4)$, $x^2{-}3x{+}2$; $P=\left[\begin{smallmatrix}1&1\\0&1\end{smallmatrix}\right]$, $D=\mathrm{diag}(2,3)$, $x^2{+}x$. (d) $T(x,y,z)=(2x,y,-2z)$, $f(x)=x^3{-}2x{+}1$: compute $f(T)$.}
\bi{EV2}{Eigenvalues and a basis of each eigenspace; is it diagonalizable, and if so find $P,D$ with $A=PDP^{-1}$: $\begin{pmatrix}2&0\\0&5\end{pmatrix}$, $\begin{pmatrix}3&1\\0&3\end{pmatrix}$, $\mathrm{diag}(1,4,-2)$, $\begin{pmatrix}2&1\\1&2\end{pmatrix}$, $\begin{pmatrix}0&1\\-2&3\end{pmatrix}$, $\begin{pmatrix}1&0&1\\0&2&0\\0&0&3\end{pmatrix}$, $\begin{pmatrix}4&0&0\\0&1&0\\0&0&2\end{pmatrix}$.}
\bi{EV7}{$T,S:\mathbb C^2\to\mathbb C^2$, $[T]=\begin{pmatrix}2&i\\-i&2\end{pmatrix}$, $[S]=\begin{pmatrix}\frac52&\frac i2\\\frac i2&\frac32\end{pmatrix}$. (a) Prove $T$ diagonalizable, $S$ not. (b) Eigenbasis $\beta$ for $T$. (c) $[T]_\beta^\beta$ and $[I]_\beta^\alpha$. (d) $[I]_\alpha^\beta$ by Gauss--Jordan. (e) $[p(T)]$ for $p=x^2{-}3x{+}1$. (f) $p(T)$ explicitly.}
\bi{CB1}{$\alpha,\beta,\gamma$ bases of $V$. Prove: $[v]_\beta=[I]_\alpha^\beta[v]_\alpha$; $[I]_\beta^\alpha=([I]_\alpha^\beta)^{-1}$; $[I]_\alpha^\gamma=[I]_\beta^\gamma[I]_\alpha^\beta$; $[I]_\beta^\alpha[I]_\alpha^\beta=I$ (via $II=I$); $[T]_\beta^\beta=[I]_\alpha^\beta[T]_\alpha^\alpha[I]_\beta^\alpha$, hence different matrices of one $T$ are similar.}
\bi{CB2}{(a) $\alpha=\{(1,0),(0,1)\}$, $\beta=\{(1,1),(1,-1)\}$: $[I]_\beta^\alpha,[I]_\alpha^\beta$, $[v]_\beta$ for $v=(4,2)$. (b) $\alpha$ standard of $\mathbb C^3$, $\beta=\{(1,0,i),(0,1,1),(1,i,0)\}$: $[I]$'s and $[v]_\beta$, $v=(2,3,1)$. (c) $T(x,y)=(x{+}2iy,3x{+}y)$ on $\mathbb C^2$, $\beta=\{(1,i),(1,-i)\}$: $[T]_\alpha^\alpha$, $[I]$'s, $[T]_\beta^\beta$ by change of basis. (d) $T(p)=p'+p$ on $\mathbf P_2$, $\alpha=\{1,x,x^2\}$, $\beta=\{1,1{+}x,(1{+}x)^2\}$: $[I]$'s, $[T]_\alpha^\alpha$, $[T]_\beta^\beta$.}

\sub{11. Orthogonality}
\bi{OR1}{\textit{(axioms)} Prove $\langle x|y\rangle=\sum x_iy_i$ on $\mathbb R^n$ is additive and homogeneous in the first slot, symmetric, positive-definite; hence linear in each slot.}
\bi{OR2}{$W\le\mathbb R^n$, $W^\perp=\{v:\langle v|w\rangle=0\ \forall w\in W\}$. Prove $W^\perp\le\mathbb R^n$. \textit{(challenge)} Prove $W\cap W^\perp=\{0\}$ and $W+W^\perp=\mathbb R^n$, i.e.\ $W\oplus W^\perp=\mathbb R^n$.}
\bi{OR3}{$P_v(x)=\frac{\langle v|x\rangle}{\langle v|v\rangle}v$ on $\mathbb R^n$. Prove: $P_v$ linear; $\mathrm{image}=\mathrm{span}(v)$, $\ker=\mathrm{span}(v)^\perp$; $P_v\circ P_v=P_v$. Reflection $R_v=2P_v-I$: linear, fixes $\mathrm{span}(v)$, negates $\mathrm{span}(v)^\perp$, $R_v\circ R_v=I$. Rotation $A_\theta$ preserves $\langle\cdot|\cdot\rangle$.}
\bi{OR4}{\textit{(fluency)} $P_{(1,2,1)}(1,-3,5)$; $P_{(2,-1,2)}(3,0,-1)$; $P_{(1,1,0)}(4,-2,3)$; $P_{(0,1,1)}(2,5,-1)$; rotation $R_{\pi/2}(3,-1)$; $R_\pi(2,5)$; $R_{\pi/3}(4,0)$; $R_{\pi/4}(2,-2)$; $R_{3\pi/2}(-1,3)$; $R_{5\pi/6}(6,0)$.}

\sub{12. Least Squares}
No new items. None of the 25 sources assesses least squares, Gram--Schmidt, QR, LU, SVD or orthogonal diagonalization; adding material would violate the efficiency criterion.

\sub{13. Other (recall and workflow)}
\bi{OT1}{\textit{Part-A recall list (one mark each).} Define: field; vector space; $\mathbb C$ operations; $\mathcal F(\mathbb R)$ operations; matrix operations; subspace; closure characterization; linear combination; span; sum of sets; LD; LI; redundant vector; linear system; REF; RREF; pivot; free/basic variable; homogeneous/non-homogeneous/inconsistent/independent-consistent/dependent-consistent system; Extension Lemma; basis; Dimension Bound Theorem; dimension; finite/infinite-dimensional; linear transformation; coordinate vector/matrix; matrix-defined map; fundamental theorem of matrices; kernel; nullspace; image; range; rank; nullity; Rank--Nullity; injective; surjective; bijective; inverse; isomorphism; invertible matrix; multilinear; alternating; determinant; Cramer's rule; diagonal; similar; diagonalizable; eigenvalue/eigenvector/eigenspace (matrix and map); change-of-basis matrix.}
